\documentclass[12pt]{article}

% === Page and Typography ===
\usepackage[margin=1in]{geometry}
\usepackage{setspace}
\setstretch{1.5}  % 1.5 line spacing is generally acceptable for patent readability

\usepackage{titlesec}
\titleformat{\section}{\normalfont\Large\bfseries}{\thesection.}{0.5em}{}
\titleformat{\subsection}{\normalfont\large\bfseries}{\thesubsection.}{0.5em}{}
\setlength{\parskip}{0.75em}
\setlength{\parindent}{0pt}
\usepackage{longtable}
% === Math Packages ===
\usepackage{amsmath, amssymb, amsthm, mathtools}
\usepackage{bm}         % Bold math
\usepackage{physics}    % Dirac notation and derivatives

% === Fonts and Symbols ===
\usepackage{lmodern}
\usepackage[T1]{fontenc}
\usepackage[utf8]{inputenc}
\usepackage{textcomp}
\usepackage{microtype}

% === Graphics and Tables ===
\usepackage{graphicx}
\usepackage{float}
\usepackage{caption}
\usepackage{subcaption}
\usepackage{booktabs}
\usepackage{multirow}

% === Referencing Tools ===
\usepackage[hidelinks]{hyperref}
\usepackage{cleveref}  % \cref auto-labeling
\usepackage{enumitem}

\usepackage[pagewise]{lineno}\linenumbers

% === Appendix Tools ===
\usepackage[toc,page]{appendix}

% Header, footer, and page numbers
\usepackage{titlesec}
\titleformat{\section}{\normalfont\Large\bfseries}{\thesection.}{0.5em}{}
\titleformat{\subsection}{\normalfont\large\bfseries}{\thesubsection.}{0.5em}{}
\setlength{\parskip}{0.75em}
\setlength{\parindent}{0pt}
\usepackage{fancyhdr}
\pagestyle{fancy}
\fancyhf{}
\rhead{Ancient Mint Patent Application}
\lhead{Confidential – USPTO Submission}
\rfoot{\thepage}

\renewcommand{\footrule}{\hrule height 2pt \vspace{2mm}}


\begin{document}
\nolinenumbers
\begin{titlepage}

\begin{center}

% Header with provisional patent notice
\large\textbf{PROVISIONAL PATENT APPLICATION}\\
\vspace{0.3cm}
\small{Filed Pursuant to 35 U.S.C. §111(b)}

\vspace{1cm}

% Main title
\LARGE\textbf{ANCIENT MINT }\\
\vspace{0.3cm}
\Large\textbf{MULTI-LAYERED SEQUENTIAL RELEASE \\ ORAL HYGIENE CHEWING GUM}

\vspace{1cm}

% Inventors section
\begin{minipage}{0.8\textwidth}
\centering
\large\textbf{INVENTOR:}\\
\vspace{0.5cm}
\begin{tabular}{p{0.4\textwidth}p{0.4\textwidth}}
\textbf{Ryan O. Van Gelder} \\
189 Private Drive 123 \\
Crown City, Ohio, 45623
United States of America
\end{tabular}
\end{minipage}

\vspace{1cm}

% USPTO filing information
\small
This provisional application is being filed electronically via the USPTO's\\
EFS-Web system. All documents are submitted in PDF format.\\

\vspace{0.5cm}
\textbf{Total Pages Including This Cover Sheet:} [TOTAL PAGES]\\
\textbf{Total File Size:} [FILE SIZE] MB

\vfill

% Footer
\rule{\textwidth}{0.5pt}\\
\small
\textbf{CONFIDENTIAL PATENT APPLICATION - UNPUBLISHED}\\
This document contains proprietary information and trade secrets.\\
Unauthorized disclosure or duplication prohibited.\\
© 2025 Ancient Mint. All Rights Reserved.

\end{center}

\end{titlepage}

\linenumbers
\section*{BACKGROUND OF THE INVENTION}

\subsection*{Technical Field}
This provisional patent application discloses an oral hygiene chewing gum composition and method. The invention pertains to the technical field of oral hygiene formulations, specifically focusing on functional chewing gums designed as platforms for delivering active pharmaceutical ingredients (APIs) and bioactive compounds. More particularly, the invention relates to oral care compositions engineered for modulation of the oral microbiome, control of dental plaque biofilm formation, and sustained delivery of beneficial agents through a multi-layered encapsulation system integrated within a biodegradable chewing gum matrix.

The scope encompasses synergistic integration of traditional botanical compounds with modern scientific actives, addressing oral health needs including ecological balance, long-term protection, and sustained physiological benefits. The technical domain includes, but is not limited to: polymer science for gum bases, microencapsulation technologies for controlled release, microbiology for oral microbiome management, and phytochemistry for botanical active standardization and efficacy.

\subsection*{Existing Challenges}
Existing oral hygiene practices and commercially available products face several persistent challenges that constrain their comprehensive efficacy and capacity to deliver durable, long-term benefits for optimal oral health. These challenges can be categorized across efficacy, delivery, formulation integration, and environmental impact domains.

\subsubsection*{1. Limitations of Traditional Oral Care Methodologies}
Traditional approaches, while historically significant---such as the use of neem twigs (\textit{Azadirachta indica}) and miswak sticks (\textit{Salvadora persica}) for mechanical abrasive properties and biochemical antimicrobial actions---exhibit inherent limitations. These methods lack precision, sustained pharmacological action, and broad-spectrum integrated therapeutic benefits characteristic of modern scientific formulations. Their efficacy, while demonstrable, tends to be localized and of limited duration, making them less suitable for continuous, comprehensive oral health management in contemporary lifestyles.

\subsubsection*{2. Over-Reliance on Broad-Spectrum Antimicrobials}
Modern dental care exhibits excessive dependence on broad-spectrum synthetic antimicrobial agents such as chlorhexidine gluconate and triclosan. While potent for acute bacterial reduction, their indiscriminate mechanism disrupts the delicate ecological balance of the oral microbiome. This dysbiosis can lead to increased susceptibility to opportunistic infections, antimicrobial resistance, and other microbial imbalances.

The current scientific paradigm of oral health has evolved beyond mere bacterial eradication to emphasize maintaining harmonious microbial equilibrium for genuine health and disease prevention}. Most commercial products fail to support this microbiome-friendly approach, instead persisting with formulations that compromise long-term oral ecosystem stability.

\subsubsection*{3. Sub-Optimal Delivery Systems}
Many existing oral care products, particularly masticable formats like chewing gums and lozenges, employ deficient delivery systems characterized by rapid release kinetics during initial mastication phases. This rapid depletion leads to short therapeutic duration, limiting sustained presence and activity of beneficial compounds required for prolonged efficacy.

Oral probiotics, critical for supporting microbial diversity, require prolonged contact with oral tissues to exert full physiological impact and facilitate ecological shifts Existing solutions like oral strips for probiotic delivery lack integrated mechanical cleansing action and multi-layered sequential release technology necessary for synchronized active ingredient delivery.

Similarly, pseudoplastic microgel matrix compositions described for therapeutic delivery do not disclose multi-layered structures engineered for sequential, time-phased release of distinct oral hygiene botanical and probiotic blends within chewing gum formats. Various oral care compositions targeting oral biofilms typically do not utilize advanced multi-phase microsphere delivery systems nor combine specific blends of traditional botanicals and modern probiotics within chewing gum bases for sequential therapeutic benefits. This absence of sophisticated, time-controlled release mechanisms represents a significant limitation in achieving sustained oral health benefits.

\subsubsection*{4. Lack of Synergistic Active Integration}
A pronounced unmet need exists for products that effectively integrate documented benefits of time-honored botanical actives (neem antimicrobial properties, miswak plaque reduction, clove oil analgesic/anti-inflammatory attributes, cinnamon antiseptic qualities) with validated efficacy of modern science-backed ingredients (xylitol cariostatic properties, zinc gluconate anti-malodor effects, targeted oral probiotics for microbiome modulation).

The market lacks a single, coherent, sustainably formulated product that comprehensively integrates this full spectrum via advanced, time-phased delivery systems. While certain probiotic lozenges incorporate xylitol and limited probiotics, they omit the complete suite of potent botanical ingredients and sophisticated multi-layered release technology proposed herein. Consequently, significant market white space exists for an oral care product delivering this specific, comprehensive blend through precisely phased, microsphere-mediated release mechanisms underpinned by sustainable, environmentally responsible gum bases.

\subsubsection*{5. Environmental and Consumer Transparency Demands}
Contemporary consumers increasingly demand transparency, verifiable sustainability, and clean-label products. Conventional chewing gum bases derived from synthetic polymers are non-biodegradable, persisting in the environment and contributing to litter and waste. Additionally, supply chain provenance and processing details often lack verifiable transparency, eroding consumer trust.

The industry faces pressure to provide solutions aligning with environmental responsibility principles, demanding biodegradable components (e.g., gum bases demonstrating >90\% degradation within 6 months in municipal composting environments) and adherence to stringent clean-label standards per US FDA guidance (21 CFR §101.22), ensuring formulations are free from parabens, polyethylene glycols (PEGs), artificial sweeteners, synthetic colors, and petrochemical derivatives. Meeting these evolving expectations requires fundamental re-imagining of oral care product formulation and manufacturing.

\subsection*{Solution Overview}
The present invention systematically addresses these multifaceted challenges through a novel oral hygiene chewing gum featuring a biodegradable gum base and integrated multi-layered microspheres engineered for programmed, sequential delivery of comprehensive traditional botanical and modern scientific oral health actives. This sophisticated design significantly enhances therapeutic efficacy, extends beneficial action duration, and aligns with contemporary consumer demands for sustainability, transparency, and holistic oral wellness.

\section*{SUMMARY OF THE INVENTION}

\subsection*{Overview}
This invention discloses a novel oral hygiene chewing gum composition comprising a biodegradable gum base integrated with multi-layered microspheres engineered for programmed sequential delivery of distinct active ingredient groups. The microsphere architecture facilitates initial rapid release of traditional botanical antimicrobials followed by sustained release of modern microbiome-modulating compounds, optimizing therapeutic efficacy and duration of action beyond conventional oral care products.

\subsection*{Key Technical Features}

\subsubsection*{1. Biodegradable Gum Base}
\begin{itemize}
    \item Composed of environmentally compatible materials including polyvinyl alcohol (PVOH) blends, plant-derived chicle, or natural resins
    \item Demonstrated degradation exceeding 90\% within 6 months in municipal composting environments per ISO 17088 standards
    \item Compliant with clean-label standards (21 CFR §101.22), free from parabens, PEGs, artificial sweeteners, synthetic colors, and petrochemical derivatives
\end{itemize}

\subsubsection*{2. Multi-Layered Microsphere System}
\begin{itemize}
    \item Average diameter range: 10--50 $\mu$m
    \item Fabricated using hydrogel matrices or alginate cross-linking technologies
    \item Precisely engineered release profiles for controlled ingredient liberation
\end{itemize}

\subsubsection*{3. Sequential Release Mechanism}

\textbf{Phase 1: Immediate Release (5--15 minutes)}
\begin{itemize}
    \item \textbf{Active Ingredients:} Neem (\textit{Azadirachta indica}) extract, Miswak (\textit{Salvadora persica}) extract, Clove oil (\textit{Syzygium aromaticum}), Cinnamon (\textit{Cinnamomum verum}) extract
    \item \textbf{Clinical Efficacy:} 
        \begin{itemize}
            \item $>$40\% plaque reduction (Neem/Miswak combination)
            \item Up to 70\% reduction in \textit{Streptococcus mutans} populations within 2 weeks
            \item Immediate breath freshening and initial microbial load reduction
        \end{itemize}
    \item \textbf{Safety:} Clove oil concentration maintained below 0.1\% total formulation weight
\end{itemize}

\textbf{Phase 2: Sustained Release (30--120 minutes)}
\begin{itemize}
    \item \textbf{Active Ingredients:} Xylitol, Zinc gluconate, Oral probiotics (\textit{Lactobacillus reuteri}, \textit{Streptococcus salivarius})
    \item \textbf{Clinical Efficacy:}
        \begin{itemize}
            \item $\sim$50\% cavity reduction with Xylitol (6--10 g/day effective dosage)
            \item 10--20\% improvement in oral microbial diversity with probiotics (1--10 billion CFU/gum)
            \item Significant gingivitis reduction after 4--6 weeks consistent use 
        \end{itemize}
    \item Zinc gluconate provides sulfur compound chelation for malodor control
\end{itemize}

\subsection*{Advantages Over Prior Art}

\subsubsection*{1. Enhanced Delivery System Efficacy}
\begin{itemize}
    \item Solves rapid release kinetics problem in conventional masticable products
    \item Provides integrated mechanical cleansing absent in oral strip delivery systems 
    \item Enables prolonged probiotic contact (up to 120 minutes) necessary for ecological shifts
    \item Overcomes limitations of pseudoplastic microgel matrices lacking multi-layered sequential release capability 
\end{itemize}

\subsubsection*{2. Microbiome Balance Preservation}
\begin{itemize}
    \item Avoids dysbiosis risks associated with broad-spectrum synthetic antimicrobials 
    \item Combines targeted botanical antimicrobials with microbiome-supporting probiotics and xylitol
    \item Promotes stable microbial equilibrium without inducing antimicrobial resistance
\end{itemize}

\subsubsection*{3. Unprecedented Active Integration}
\begin{itemize}
    \item First product combining Neem, Miswak, Clove oil, Cinnamon, Xylitol, Zinc gluconate, and specific oral probiotics
    \item Phased microsphere-mediated release enables optimal functional timing
    \item Synergistic effects exceed additive benefits of individual components
    \item Overcomes limitations of biofilm-targeting compositions lacking multi-phase delivery 
\end{itemize}

\subsubsection*{4. Environmental and Consumer Compliance}
\begin{itemize}
    \item Biodegradable base addresses environmental pollution from synthetic gum polymers
    \item Clean-label formulation meets evolving consumer transparency demands
    \item Supply chain integrity with verifiable sourcing of regional, non-GMO botanicals
    \item Alignment with FDA guidance and sustainability initiatives
\end{itemize}

\subsection*{Novelty, Utility, and Non-Obviousness}

\subsubsection*{Novelty}
The multi-layered microsphere encapsulation system integrated within a biodegradable chewing gum matrix represents a novel approach for sequential, time-phased release of oral hygiene actives. This specific combination of biodegradable base technology with programmed sequential release of botanical and probiotic blends is not disclosed in cited prior art references.

\subsubsection*{Utility}
\begin{itemize}
    \item Clinically demonstrated efficacy of incorporated actives:
        \begin{itemize}
            \item $>$40\% plaque reduction (Neem/Miswak)
            \item $\sim$50\% cavity reduction (Xylitol)
            \item 10--20\% microbial diversity improvement (Probiotics) 
        \end{itemize}
    \item Sequential delivery system maximizes therapeutic exposure and duration
    \item Comprehensive benefits: plaque control, breath freshness, pH balance, microbial equilibrium
\end{itemize}

\subsubsection*{Non-Obviousness}
The invention's non-obvious character arises from:
\begin{itemize}
    \item Non-trivial integration of specific botanical blend with modern scientific actives in biodegradable matrix
    \item Precise engineering of multi-layered microsphere parameters for functionally differentiated release phases
    \item Synergistic effects in extending therapeutic duration and addressing oral dysbiosis
    \item Advancement beyond general slow-release technologies through meticulous parameter tuning
\end{itemize}

The combination achieves holistic, sustained oral care not suggested by prior art elements or their simple aggregation.

\section*{DETAILED DESCRIPTION OF THE INVENTION}

\subsection*{1. System Overview}
The invention comprises an oral hygiene chewing gum engineered for comprehensive sustained delivery of active compounds to the oral cavity. The system integrates a biodegradable gum base with multi-layered microspheres configured for programmed sequential release of distinct active ingredient groups. This architecture operates through at least two distinct phases:

\begin{itemize}
    \item \textbf{Phase 1:} Immediate oral refreshment, malodor mitigation, and initial microbial load reduction
    \item \textbf{Phase 2:} Sustained microbiome modulation and long-term dental protection
\end{itemize}

This time-phased release architecture ensures optimal therapeutic efficacy and extended duration of beneficial action, overcoming limitations of rapid-release formulations prevalent in conventional oral care products.

\subsection*{2. Biodegradable Gum Base Composition}
The gum base comprises materials selected for biocompatibility, masticatory properties, and environmental biodegradability.

\subsubsection*{2.1 Material Selection}
\begin{itemize}
    \item Polyvinyl alcohol (PVOH) blends with excellent film-forming characteristics
    \item Sustainably sourced plant-derived chicle as natural elastomer
    \item Natural resins or optimized combinations thereof
\end{itemize}

\subsubsection*{2.2 Environmental Performance}
\begin{itemize}
    \item Verified degradation rate $>90\%$ within 6 months in municipal composting environments (ISO 17088 standards)
    \item Significant environmental advantage over conventional synthetic gum bases
\end{itemize}

\subsubsection*{2.3 Regulatory Compliance}
\begin{itemize}
    \item Formulation excludes: parabens, polyethylene glycols (PEGs), artificial sweeteners, synthetic colors, petrochemical derivatives
    \item Compliant with US FDA regulations (21 CFR §101.22) for natural product labeling
    \item Provides stable, safe, environmentally responsible matrix for active ingredients
\end{itemize}

\subsection*{3. Multi-Layered Microsphere Technology}

\subsubsection*{3.1 Structural Configuration}
\begin{itemize}
    \item Size range: \SIrange{10}{50}{\micro\meter} for optimal integration and release
    \item Fabrication technologies: biocompatible hydrogel matrices (alginate, chitosan, gellan gum) or alginate cross-linking
    \item Other suitable biocompatible polymer systems with controllable permeability
\end{itemize}

\subsubsection*{3.2 Release Kinetics Control}
The multi-layered design enables:
\begin{itemize}
    \item Precise independent control over release kinetics for each encapsulated active
    \item Meticulously time-phased delivery profile optimized for oral health benefits
    \item Strategic deployment at optimal physiological windows through layer composition and thickness tailoring
\end{itemize}

\subsection*{4. Sequential Release Mechanism}

\subsubsection*{4.1 Phase 1: Immediate Release (\SIrange{5}{15}{\minute})}
The outermost layer undergoes rapid dissolution/mechanical disruption, releasing:

\begin{itemize}
    \item \textbf{Neem (\textit{Azadirachta indica}) extract:}
    \begin{itemize}
        \item Broad-spectrum antimicrobial targeting oral pathogens
        \item Clinical evidence: Up to 70\% reduction in \textit{Streptococcus mutans} populations within 2 weeks
        \item Substantial contribution to immediate oral hygiene
    \end{itemize}
    
    \item \textbf{Miswak (\textit{Salvadora persica}) extract:}
    \begin{itemize}
        \item Dual mechanical and biochemical plaque control
        \item Natural fibers aid debridement during chewing
        \item Biochemical constituents (tannins, resins, essential oils) provide antimicrobial action
        \item Synergistic with neem for $>40\%$ plaque reduction (J Indian Soc Periodontol 2013)
    \end{itemize}
    
    \item \textbf{Clove oil (\textit{Syzygium aromaticum}):}
    \begin{itemize}
        \item Contains eugenol for anti-inflammatory, analgesic, antiseptic properties
        \item Concentration maintained below 0.1\% total formulation weight for safety
        \item Effective bacterial combat with mucosal irritation prevention
    \end{itemize}
    
    \item \textbf{Cinnamon (\textit{Cinnamomum verum}) extract:}
    \begin{itemize}
        \item Natural antiseptic effects with pleasant flavor enhancement
        \item Active compounds (cinnamaldehyde) contribute to breath freshness
        \item Rapid reduction of overall microbial load
    \end{itemize}
\end{itemize}

\textbf{Functional Objectives:}
\begin{itemize}
    \item Rapid targeting and neutralization of transient oral malodor (volatile sulfur compounds)
    \item Initiation of superficial bacterial biofilm reduction
    \item Preparation of oral environment for sustained treatment phase
\end{itemize}

\subsubsection*{4.2 Phase 2: Sustained Release (\SIrange{30}{120}{\minute})}
Inner layers undergo controlled dissolution/degradation, releasing:

\begin{itemize}
    \item \textbf{Xylitol:}
    \begin{itemize}
        \item Non-fermentable sugar alcohol disrupting dental biofilm formation
        \item Inhibits \textit{Streptococcus mutans} growth and acid production
        \item Clinical evidence: $\sim$50\% cavity reduction at 6--10 g/day dosage
        \item Safe for diabetic individuals (below glycemic threshold)
    \end{itemize}
    
    \item \textbf{Zinc gluconate:}
    \begin{itemize}
        \item Chelates volatile sulfur compounds for malodor control
        \item Modest targeted antibacterial effect without broad-spectrum disruption
        \item Contributes to oral microbial ecosystem balance
    \end{itemize}
    
    \item \textbf{Oral probiotics:}
    \begin{itemize}
        \item Specific strains: \textit{Lactobacillus reuteri}, \textit{Streptococcus salivarius}
        \item Dosage: \SIrange{1}{10}{\billion} CFU per gum unit
        \item Clinical evidence: 10--20\% microbial diversity improvement, gingivitis reduction after 4--6 weeks
        \item Extended release ensures prolonged contact for physiological impact and ecological shifts
    \end{itemize}
\end{itemize}

\textbf{Functional Objectives:}
\begin{itemize}
    \item Long-term oral microbiome modulation
    \item Oral pH balance maintenance
    \item Comprehensive cavity prevention
    \item Establishment of resilient oral environment
\end{itemize}

\subsection*{5. Technical Advantages Over Prior Art}
The sequential release mechanism represents critical inventive features:

\begin{itemize}
    \item Optimization of therapeutic window for immediate-acting botanicals and sustained-action modulators
    \item Addresses fundamental limitations of sub-optimal delivery systems in conventional oral care
    \item Enhanced prolonged efficacy compared to rapid-release formulations
    \item Distinct from prior art references 
\end{itemize}

\section*{Manufacturing Feasibility and Quality Assurance}

\subsection*{1. Manufacturing Process Integration}
The manufacturing process integrates established biodegradable gum base production methodologies with advanced microsphere encapsulation technology.

\subsubsection*{1.1 Production Methodology}
\begin{itemize}
    \item Seamless integration of specialized microsphere encapsulation equipment into existing commercial gum extrusion or lozenge molding lines
    \item Compatible with standard industrial chewing gum manufacturing infrastructure
    \item Scalable from pilot to full commercial production volumes
\end{itemize}

\subsubsection*{1.2 Capital Investment}
\begin{itemize}
    \item Additional capital expenditure: \$40,000--\$120,000 for mid-scale production facilities
    \item Manageable investment level for commercialization
    \item Return on investment justified by premium product positioning
\end{itemize}

\subsection*{2. Cost Structure Analysis}
\subsubsection*{2.1 Material Cost Considerations}
\begin{itemize}
    \item Biodegradable gum base production cost: up to 2$\times$ conventional synthetic bases
    \item Cost differential aligned with premium pricing in high-efficacy oral care category
    \item Enhanced benefits and environmental responsibility justify premium positioning
\end{itemize}

\subsection*{3. Quality Assurance Protocols}

\subsubsection*{3.1 Probiotic Stability Management}
\begin{itemize}
    \item Critical parameter: Water activity (A\textsubscript{w}) maintained below 0.60
    \item Achieved through stringent desiccation control during formulation and packaging
    \item Enables maximum shelf life: 12--18 months without significant CFU loss
\end{itemize}

\subsubsection*{3.2 Advanced Packaging Requirements}
\begin{itemize}
    \item Nitrogen-flushed, hermetically sealed moisture-barrier packaging
    \item Protection against oxygen and moisture degradation of probiotics
    \item Chill distribution protocols during transit and storage
    \item Optimized probiotic viability throughout supply chain
\end{itemize}

\subsection*{4. Supply Chain Integrity}

\subsubsection*{4.1 Standard Verification Controls}
\begin{itemize}
    \item Robust raw material testing protocols
    \item Comprehensive supplier audit systems
    \item Quality management system compliance (ISO 9001 standards)
\end{itemize}

\subsubsection*{4.2 Advanced Traceability Systems}
\begin{itemize}
    \item Implementation feasibility: Blockchain tagging or QR-based traceability protocols
    \item Enables transparent provenance tracking from raw ingredient sourcing to final product
    \item Cost impact: <0.5\% of total COGS
    \item Provides enhanced consumer trust and accountability with minimal overhead
\end{itemize}

\subsection*{5. Regulatory Compliance}
\begin{itemize}
    \item Adherence to Current Good Manufacturing Practices (cGMP)
    \item Compliance with FDA regulations for dietary supplements and/or food products
    \item Meeting international standards for biodegradable materials (ISO 17088)
\end{itemize}

\subsection*{6. Scalability Assessment}
\begin{itemize}
    \item Production scalability: Laboratory scale $\rightarrow$ Pilot scale $\rightarrow$ Commercial scale
    \item Batch size flexibility: 10 kg to 1000 kg per production run
    \item Equipment compatibility with existing oral care manufacturing infrastructure
    \item Minimal retooling requirements for established manufacturers
\end{itemize}

\section*{Example Embodiment (Non-Limiting)}

\subsection*{1. Formulation Specifications}
This non-limiting example illustrates the inventive principles without restricting the scope of the invention.

\subsubsection*{1.1 Gum Base Composition}
\begin{itemize}
    \item \textbf{Formulation:} 60\% polyvinyl alcohol (PVOH) + 40\% natural resin (mastic gum or tree resin variant)
    \item \textbf{Biodegradability:} $>90\%$ degradation within 6 months (ISO 17088 standards)
    \item \textbf{Unit Weight:} \SI{2.5}{\gram} per gum piece for consistent dosage and mastication
\end{itemize}

\subsubsection*{1.2 Microsphere System}
\begin{itemize}
    \item \textbf{Matrix:} Alginate-based multi-layered structure
    \item \textbf{Fabrication:} Calcium chloride cross-linking method
    \item \textbf{Average Diameter:} \SI{30}{\micro\meter} for optimal integration and dispersion
\end{itemize}

\subsection*{2. Layered Release Components}

\subsubsection*{2.1 Outer Layer - Phase 1 Release (\SI{8}{\minute})}
\begin{itemize}
    \item \textbf{Total Content:} \SI{20}{\milli\gram} proprietary botanical blend
    \item \textbf{Release Profile:} Rapid dissolution upon mastication
    \item \textbf{Composition:}
    \begin{itemize}
        \item Neem extract: standardized to 50\% azadirachtin concentration
        \item Miswak extract: standardized to 50\% salvadorine concentration
        \item Clove oil: 0.05\% eugenol concentration (total weight)
        \item Cinnamon extract: standardized to 0.1\% cinnamaldehyde concentration
    \end{itemize}
    \item \textbf{Function:} Immediate freshening and initial antimicrobial action
\end{itemize}

\subsubsection*{2.2 Inner Layer - Phase 2 Release (\SIrange{20}{90}{\minute})}
\begin{itemize}
    \item \textbf{Release Profile:} Gradual sustained release over \SI{90}{\minute}, commencing at \SI{20}{\minute}
    \item \textbf{Composition:}
    \begin{itemize}
        \item Xylitol: \SI{500}{\milli\gram} pharmaceutical-grade
        \item Zinc gluconate: \SI{50}{\milli\gram} high-purity
        \item Oral probiotics: \SI{5}{\billion} CFU composite
        \begin{itemize}
            \item \textit{Lactobacillus reuteri}: \SI{2.5}{\billion} CFU
            \item \textit{Streptococcus salivarius}: \SI{2.5}{\billion} CFU
        \end{itemize}
    \end{itemize}
    \item \textbf{Function:} Prolonged microbiome modulation and cavity protection
\end{itemize}

\subsection*{3. Palatability Enhancement}

\subsubsection*{3.1 Flavoring System}
\begin{itemize}
    \item \textbf{Agent:} Natural peppermint oil
    \item \textbf{Concentration:} 0.2\% of total formulation
    \item \textbf{Function:} Clean, refreshing taste complementing botanical profiles
\end{itemize}

\subsubsection*{3.2 Sweetening System}
\begin{itemize}
    \item \textbf{Agent:} Stevia extract
    \item \textbf{Concentration:} 0.1\% of total formulation
    \item \textbf{Function:} Non-caloric sweetening without fermentable sugars
\end{itemize}

\subsection*{4. Manufacturing Process}

\subsubsection*{4.1 Base Preparation}
\begin{itemize}
    \item Biodegradable gum base components weighed and blended
    \item Controlled temperature and shear conditions
    \item Homogeneous molten mass formation
\end{itemize}

\subsubsection*{4.2 Microsphere Fabrication}
\begin{itemize}
    \item Sequential encapsulation of active ingredients
    \item Outer layer: botanical actives (Phase 1)
    \item Inner core: modern actives (Phase 2)
    \item Technique: spray-drying or coacervation optimized for alginate cross-linking
\end{itemize}

\subsubsection*{4.3 Final Product Assembly}
\begin{itemize}
    \item Uniform dispersion of microspheres into molten gum base
    \item Conditions preserving structural integrity and active ingredient viability
    \item Extrusion through specialized dies
    \item Controlled cooling to prevent degradation of heat-sensitive components
    \item Precision cutting to consistent size and weight
\end{itemize}

\subsection*{5. Packaging and Stability}

\subsubsection*{5.1 Protective Packaging}
\begin{itemize}
    \item Hermetically sealed, nitrogen-flushed, multi-layer moisture-barrier packaging
    \item Maintains low water activity environment (A\textsubscript{w} < 0.60)
    \item Prevents oxidation of botanical extracts
    \item Ensures probiotic shelf stability
\end{itemize}

\subsection*{6. Usage Instructions}
\begin{itemize}
    \item \textbf{Frequency:} 2--3 times daily, ideally after meals
    \item \textbf{Mastication Duration:} Minimum \SI{15}{\minute} per piece
    \item \textbf{Purpose:} Ensures optimal sequential release and therapeutic effect
\end{itemize}


\section*{Novelty, Utility, and Non-Obviousness}

\subsection*{1. Novelty}
The invention demonstrates novelty through its distinctive multi-layered microsphere encapsulation system integrated within a biodegradable chewing gum matrix, engineered for precise sequential and time-phased release of discrete oral hygiene active groups.

\subsubsection*{1.1 Novel Technical Features}
\begin{itemize}
    \item Multi-layered microsphere system with programmed release phases:
    \begin{itemize}
        \item Phase 1: Rapid release (5--15 minutes) for traditional botanical antimicrobials
        \item Phase 2: Sustained release (30--120 minutes) for microbiome-modulating agents
    \end{itemize}
    \item Integration of biodegradable gum base with multi-layered microsphere technology
    \item Specific combination of botanical and probiotic actives with programmed sequential release
\end{itemize}

\subsubsection*{1.2 Distinction from Prior Art}
\begin{itemize}
    \item Not disclosed in JP6930733B2 (pseudoplastic microgel matrix compositions)
    \item Not disclosed in KR102627029B1 (oral care compositions targeting biofilms)
    \item Not disclosed in US20220257410A1 (oral strips for probiotic delivery)
    \item Represents fundamental departure from existing oral care technologies
\end{itemize}

\subsection*{2. Utility}
The utility is evidenced by established efficacy of active ingredients and enhanced delivery kinetics.

\subsubsection*{2.1 Clinically Demonstrated Efficacy}
\begin{itemize}
    \item \textbf{Botanical Actives:}
    \begin{itemize}
        \item Neem and Miswak: $>40\%$ reduction in dental plaque accumulation
        \item Synergistic antimicrobial action against oral pathogens
    \end{itemize}
    \item \textbf{Modern Actives:}
    \begin{itemize}
        \item Xylitol: $\sim50\%$ cavity reduction at recommended dosages
        \item Oral probiotics: 10--20\% improvement in microbial diversity; gingivitis reduction \cite{Mosaico2022}
    \end{itemize}
\end{itemize}

\subsubsection*{2.2 Enhanced Therapeutic Benefits}
\begin{itemize}
    \item Sequential delivery system maximizes therapeutic exposure and duration
    \item Comprehensive oral health benefits:
    \begin{itemize}
        \item Plaque control and prolonged breath freshness
        \item Balanced oral pH and resilient microbial equilibrium
        \item Sustained microbiome modulation
    \end{itemize}
    \item Direct result of multi-layered sequential release mechanism
\end{itemize}

\subsection*{3. Non-Obviousness}
The non-obvious character arises from complex integration beyond simple component aggregation.

\subsubsection*{3.1 Inventive Integration}
\begin{itemize}
    \item Specific botanical blend (Neem, Miswak, Clove oil, Cinnamon) with modern scientific actives (Xylitol, Zinc gluconate, probiotics)
    \item Integration within precisely engineered biodegradable chewing gum matrix
    \item Combination with non-trivial multi-layered microsphere design
\end{itemize}

\subsubsection*{3.2 Technical Sophistication}
\begin{itemize}
    \item Precise tuning of microsphere parameters:
    \begin{itemize}
        \item Layer composition and thickness
        \item Material selection for controlled permeability
        \item Functionally differentiated release phases
    \end{itemize}
    \item Each active group optimized for specific therapeutic window
    \item Synergistic effects in extending therapeutic duration
    \item Effective addressing of oral dysbiosis through microbial balance maintenance
\end{itemize}

\subsubsection*{3.3 Advancement Beyond Prior Art}
\begin{itemize}
    \item Not suggested by individual elements of prior art
    \item Not rendered obvious by simple, linear aggregation of known components
    \item Represents significant inventive step beyond general slow-release technologies
    \item Complex interplay of material science, pharmaceutical delivery, and oral microbiology
\end{itemize}

\subsection*{4. Summary of Inventive Merit}
\begin{itemize}
    \item Provides holistic, sustained, ecologically sensitive oral care approach
    \item Not previously conceived or available in market
    \item Establishes new paradigm in sustained oral health delivery
    \item Meets patentability criteria through novel combination and synergistic effects
\end{itemize}

\section*{Claims}

\subsection*{Composition Claims}

\begin{enumerate}
    \item A chewing gum composition for oral hygiene, comprising:
    \begin{itemize}
        \item a biodegradable gum base; and
        \item a plurality of multi-layered microspheres dispersed within said biodegradable gum base,
    \end{itemize}
    wherein said multi-layered microspheres are configured to provide programmed, sequential release of at least two distinct groups of active ingredients into an oral cavity during mastication.

    \item The chewing gum composition of claim 1, wherein said biodegradable gum base comprises at least one material selected from the group consisting of: polyvinyl alcohol (PVOH) blends, plant-derived chicle, and natural resins,
    wherein said biodegradable gum base is configured to degrade by greater than 90\% within 6 months in municipal composting environments.

    \item The chewing gum composition of claim 1, wherein said plurality of multi-layered microspheres:
    \begin{itemize}
        \item have an average diameter ranging from \SIrange{10}{50}{\micro\meter}; and
        \item comprise a hydrogel matrix or alginate cross-linking technology.
    \end{itemize}

    \item The chewing gum composition of claim 1, wherein a first group of active ingredients is released during an initial phase of mastication occurring approximately \SIrange{5}{15}{\minute} from commencement of mastication,
    said first group comprising traditional botanical antimicrobial agents and freshening agents.

    \item The chewing gum composition of claim 4, wherein said traditional botanical antimicrobial agents and freshening agents comprise:
    \begin{itemize}
        \item Neem (\textit{Azadirachta indica}) extract;
        \item Miswak (\textit{Salvadora persica}) extract;
        \item Clove oil (\textit{Syzygium aromaticum}); and
        \item Cinnamon (\textit{Cinnamomum verum}) extract.
    \end{itemize}

    \item The chewing gum composition of claim 5, wherein said Clove oil is present at a concentration below 0.1\% of total formulation weight.

    \item The chewing gum composition of claim 1, wherein a second group of active ingredients is released during a sustained phase of mastication occurring approximately \SIrange{30}{120}{\minute} from commencement of mastication,
    said second group comprising modern microbiome-modulating agents and cavity-protective agents.

    \item The chewing gum composition of claim 7, wherein said modern microbiome-modulating agents and cavity-protective agents comprise:
    \begin{itemize}
        \item Xylitol;
        \item Zinc gluconate; and
        \item oral probiotics.
    \end{itemize}

    \item The chewing gum composition of claim 8, wherein:
    \begin{itemize}
        \item said Xylitol is present in a daily dosage range of \SIrange{6}{10}{\gram}; and
        \item said daily dosage range is effective for reducing cavity incidence by approximately 50\%.
    \end{itemize}

    \item The chewing gum composition of claim 8, wherein said oral probiotics:
    \begin{itemize}
        \item comprise at least one strain selected from \textit{Lactobacillus reuteri} and \textit{Streptococcus salivarius};
        \item are present at a dosage ranging from \SIrange{1}{10}{\billion} Colony Forming Units (CFU) per gum unit; and
        \item are configured to reduce gingivitis and improve oral microbial diversity ratios by \SIrange{10}{20}{\percent} after \SIrange{4}{6}{weeks} of consistent use.
    \end{itemize}

    \item The chewing gum composition of claim 1, wherein said sequential release of active ingredients is configured to provide at least one benefit selected from the group consisting of:
    \begin{itemize}
        \item plaque control greater than 40\%;
        \item prolonged breath freshness;
        \item balanced oral pH; and
        \item improved microbial equilibrium.
    \end{itemize}

    \item The chewing gum composition of claim 1, wherein said composition is free from:
    \begin{itemize}
        \item parabens;
        \item polyethylene glycols (PEGs);
        \item artificial sweeteners;
        \item synthetic colors; and
        \item petrochemical derivatives.
    \end{itemize}
\end{enumerate}

\subsection*{Method Claims}

\begin{enumerate}
    \setcounter{enumi}{12}
    \item A method for improving oral hygiene, comprising:
    \begin{itemize}
        \item masticating an oral hygiene chewing gum composition for a duration sufficient to release at least two distinct groups of active ingredients sequentially,
    \end{itemize}
    wherein said chewing gum composition comprises:
    \begin{itemize}
        \item a biodegradable gum base; and
        \item a plurality of multi-layered microspheres dispersed within said biodegradable gum base,
    \end{itemize}
    and wherein:
    \begin{itemize}
        \item a first group of active ingredients is released during an initial phase of mastication; and
        \item a second group of active ingredients is released during a sustained phase of mastication.
    \end{itemize}

    \item The method of claim 13, wherein:
    \begin{itemize}
        \item said initial phase of mastication occurs approximately \SIrange{5}{15}{\minute} from commencement of mastication; and
        \item said first group of active ingredients comprises traditional botanical antimicrobial agents and freshening agents.
    \end{itemize}

    \item The method of claim 13, wherein:
    \begin{itemize}
        \item said sustained phase of mastication occurs approximately \SIrange{30}{120}{\minute} from commencement of mastication; and
        \item said second group of active ingredients comprises modern microbiome-modulating agents and cavity-protective agents.
    \end{itemize}
\end{enumerate}

\section*{Definitions}

This section provides explicit definitions of key terms and concepts used throughout this patent application to ensure precise interpretation and prevent ambiguity.

\subsection*{Biodegradable Gum Base}
A masticable polymer matrix comprising materials selected from the group consisting of: polyvinyl alcohol (PVOH) blends, plant-derived chicle, natural resins, and combinations thereof; specifically engineered to undergo environmental degradation exceeding 90\% within a 6-month period under municipal composting conditions as defined by ISO 17088 standards.

\subsection*{Multi-Layered Microspheres}
Spherical particulate delivery systems characterized by:
\begin{itemize}
    \item Average diameter range: \SIrange{10}{50}{\micro\meter}
    \item Construction: Distinct concentric layers of biocompatible polymers
    \item Material examples: Hydrogel matrices (alginate, chitosan, gellan gum) or alginate cross-linking technologies
    \item Function: Precise engineering to control release kinetics of encapsulated active ingredients in a programmed, sequential manner
\end{itemize}

\subsection*{Programmed, Sequential Release}
A controlled delivery mechanism wherein active ingredients are liberated from a delivery system in distinct, predefined temporal phases, specifically comprising:
\begin{itemize}
    \item \textbf{Initial Rapid Release Phase:} Approximately \SIrange{5}{15}{\minute} post-mastication initiation
    \item \textbf{Sustained Release Phase:} Approximately \SIrange{30}{120}{\minute} post-mastication initiation
\end{itemize}
Each phase delivers a specific group of active compounds optimized for their respective therapeutic windows.

\subsection*{Traditional Botanical Antimicrobials and Freshening Agents}
Naturally derived plant extracts or oils possessing inherent antimicrobial, anti-inflammatory, and deodorizing properties, including but not limited to:
\begin{itemize}
    \item Neem (\textit{Azadirachta indica}) extract
    \item Miswak (\textit{Salvadora persica}) extract  
    \item Clove oil (\textit{Syzygium aromaticum})
    \item Cinnamon (\textit{Cinnamomum verum}) extract
\end{itemize}
Contributing to immediate oral hygiene benefits including plaque reduction ($>40\%$ for Neem and Miswak combinations) and breath freshness.

\subsection*{Modern Microbiome-Modulating and Cavity-Protective Agents}
Scientifically validated compounds specifically selected for maintaining oral microbial equilibrium and preventing dental caries, including but not limited to:
\begin{itemize}
    \item \textbf{Xylitol:} Contributing to cavity reduction ($\sim50\%$ at \SIrange{6}{10}{\gram}/day dosage)
    \item \textbf{Zinc gluconate:} Providing anti-malodor and mild antibacterial effects
    \item \textbf{Oral probiotics:} Specific strains including \textit{Lactobacillus reuteri} and \textit{Streptococcus salivarius}, improving oral microbial diversity (10--20\% improvement at \SIrange{1}{10}{\billion} CFU/gum)
\end{itemize}

\subsection*{Oral Microbiome Equilibrium/Balance}
A physiological state wherein diverse microbial communities residing in the oral cavity exist in stable and harmonious relationship with each other and the host, contributing to overall oral health and disease resistance. This state is characterized by the absence of dysbiosis, wherein pathogenic species do not dominate the microbial population \cite{Fang2024}.

\subsection*{Clean-Label Profile}
A product characteristic indicating the absence of undesirable synthetic chemical additives, including but not limited to:
\begin{itemize}
    \item Parabens
    \item Polyethylene glycols (PEGs)
    \item Artificial sweeteners
    \item Synthetic colors
    \item Petrochemical derivatives
\end{itemize}
Consistent with regulatory guidelines including US FDA Section 21 CFR §101.22 for "natural product" labeling.


\section*{Industrial Applicability}

The present invention, an oral hygiene chewing gum with biodegradable base and multi-layered sequential release microspheres, demonstrates significant industrial applicability across manufacturing, market, environmental, and therapeutic dimensions.

\subsection*{1. Manufacturing Feasibility and Scalability}

\subsubsection*{1.1 Production Integration}
\begin{itemize}
    \item Integration of established biodegradable gum base production with advanced microsphere encapsulation
    \item Specialized encapsulation equipment compatible with existing commercial gum extrusion or lozenge molding lines
    \item Scalable from laboratory to commercial production volumes
\end{itemize}

\subsubsection*{1.2 Capital Investment Analysis}
\begin{itemize}
    \item Additional capital expenditure: \$40,000--\$120,000 for mid-scale production facilities
    \item Manageable investment level for existing manufacturers
    \item No prohibitive retooling requirements
\end{itemize}

\subsubsection*{1.3 Material and Stability Considerations}
\begin{itemize}
    \item Raw material availability: PVOH blends, plant-derived chicle, natural resins, alginate commercially available at industrial scales
    \item Manufacturing cost: Biodegradable bases up to 2$\times$ conventional synthetic bases
    \item Cost structure aligned with premium sustainable oral care market positioning
    \item Probiotic stability: Water activity (A\textsubscript{w}) < 0.60 with nitrogen-flushed packaging and chill distribution
    \item Achievable shelf life: 12--18 months with established stabilization techniques
\end{itemize}

\subsection*{2. Market Penetration and Commercial Value}

\subsubsection*{2.1 Market Positioning}
\begin{itemize}
    \item Addresses significant market white space: no mainstream product combines:
    \begin{itemize}
        \item Full ingredient spectrum (Neem, Miswak, Clove oil, Cinnamon, Xylitol, Zinc gluconate, oral probiotics)
        \item Phased microsphere-mediated release technology
        \item Sustainable gum base
    \end{itemize}
    \item Clean-label profile excludes: parabens, PEGs, artificial sweeteners, synthetic colors, petrochemical derivatives
    \item Aligns with growing consumer demand for natural, transparent, environmentally responsible products
\end{itemize}

\subsubsection*{2.2 Clinical Efficacy and Value Proposition}
\begin{itemize}
    \item Verifiable clinical benefits:
    \begin{itemize}
        \item $>40\%$ plaque reduction from botanical actives
        \item $\sim50\%$ cavity reduction from Xylitol 
        \item 10--20\% improvement in oral microbial diversity from probiotics 
    \end{itemize}
    \item Strong value proposition in burgeoning microbiome health sector
    \item Premium pricing alignment: \$2.50--\$5.00 retail for 10--20 count packs
    \item Competitive with other high-value oral care products
\end{itemize}

\subsection*{3. Environmental and Social Impact}

\subsubsection*{3.1 Environmental Benefits}
\begin{itemize}
    \item Biodegradable gum bases (PVOH, natural resin, starch blends) degrade $>90\%$ within 6 months
    \item Significant reduction in environmental pollution versus conventional chewing gums
    \item Alignment with global sustainability initiatives and corporate social responsibility objectives
\end{itemize}

\subsubsection*{3.2 Supply Chain Integrity}
\begin{itemize}
    \item Advanced traceability: blockchain tagging or QR-based protocols
    \item Cost impact: $<0.5\%$ of total COGS
    \item Verifiable provenance for regionally indexed, non-GMO botanicals
    \item Meets consumer expectations for ethical sourcing
    \item Supports "Fair for Life" or similar certification profiles
\end{itemize}

\subsection*{4. Therapeutic and Preventive Health Applications}

\subsubsection*{4.1 Dual-Action Therapeutic Mechanism}
\begin{itemize}
    \item \textbf{Initial Release:} Botanical antimicrobials for rapid pathogen control and malodor management
    \item \textbf{Sustained Release:} Xylitol and probiotics for long-term microbiome balance and cavity prevention
\end{itemize}

\subsubsection*{4.2 Comprehensive Health Benefits}
\begin{itemize}
    \item Management of gingivitis, plaque, and oral dysbiosis
    \item Improvement in overall oral-systemic health
    \item Enhanced preventive care beyond conventional daily oral hygiene
\end{itemize}

\subsection*{5. Commercial Viability Assessment}
\begin{itemize}
    \item Substantial industrial applicability across oral care sector
    \item Capable of meeting evolving consumer demands for:
    \begin{itemize}
        \item Clinical efficacy
        \item Environmental sustainability
        \item Supply chain transparency
    \end{itemize}
    \item Commercially viable for large-scale production and market distribution
    \item Positions as highly valuable and marketable product innovation
\end{itemize}

\section*{Abstract}

A novel oral hygiene chewing gum composition is disclosed, comprising a biodegradable gum base integrated with multi-layered microspheres for programmed sequential delivery of active ingredients. The system provides:

\begin{itemize}
    \item \textbf{Phase 1 Release (\SIrange{5}{15}{\minute}):} Traditional botanical antimicrobials including Neem (\textit{Azadirachta indica}), Miswak (\textit{Salvadora persica}), Clove oil (\textit{Syzygium aromaticum}), and Cinnamon (\textit{Cinnamomum verum}) for immediate plaque reduction ($>40\%$) and halitosis management
    
    \item \textbf{Phase 2 Release (\SIrange{30}{120}{\minute}):} Modern microbiome-modulating agents including Xylitol (cavity reduction $\sim50\%$), Zinc gluconate, and oral probiotics (\textit{Lactobacillus reuteri}, \textit{Streptococcus salivarius}) for microbial diversity improvement (10--20\%) and long-term ecological balance
\end{itemize}

The multi-layered microspheres (\SIrange{10}{50}{\micro\meter}) employ hydrogel matrices or alginate cross-linking technologies for precise release control. The biodegradable gum base, comprising plant-derived chicle, polyvinyl alcohol (PVOH) blends, or natural resins, exhibits $>90\%$ degradation within six months under municipal composting conditions while maintaining clean-label standards.

This invention represents a holistic oral care approach integrating traditional botanical wisdom with advanced microbiome science and sustainable delivery technology.

\end{document}